购买系统还可以改变执行流程的保存方式、持久化的等待位置，以及状态文件的发布结构。下面的九个系统分别进入这三个方向。DBOS、Restate 和 LangGraph 保存程序或图的执行位置，回答重启后从哪里继续。libDSE、Jetpack 和 Bodega 分别改变依赖持久化、命令提交和线性化读取的等待路径。SlateDB、Iceberg 和 Delta Lake 则把计算进程与对象存储中的有效状态连接起来。

\begin{table}[htbp]
\centering\small
\begin{tabularx}{\linewidth}{@{}p{.17\linewidth}p{.24\linewidth}X@{}}
\toprule
演进方向 & 系统 & 当前设计要核对的对象与边界 \\
\midrule
执行流程的持久恢复 & DBOS、Restate、LangGraph & 已保存步骤、调用 journal 或图检查点决定重放位置，原购买身份连接付款成功到执行记录保存之间的窗口。 \\
延迟路径的重新安排 & libDSE、Jetpack、Bodega & libDSE 等待可恢复依赖闭包，Jetpack 为独立命令增加共识快路径，Bodega 以写确认覆盖和租约换取选定地域的本地线性化读。 \\
对象状态的持久化与发布 & SlateDB、Iceberg、Delta Lake & SlateDB 区分内存写入与对象持久完成，Iceberg 和 Delta 发布完整文件集合，并各自验证并发提交条件。 \\
\bottomrule
\end{tabularx}
\caption{三个演进方向分别改变执行历史、等待路径与有效状态的发布位置}
\label{tab:ch21-frontier-directions}
\end{table}

先沿执行流程看恢复位置。相同的付款步骤可以由应用内的库调度，也可以交给独立日志运行时，包含模型选择的流程还要保存决定与批准。三种结构都要在步骤重入时找到已经固定的业务输入，下面从最接近普通工作者的 DBOS 开始。

\subsection{DBOS 把工作流控制流程放进数据库}

普通支付工作者从 C/T 扫描责任，再根据当前状态决定下一步。持久执行进一步保存程序走到了哪里：一份工作流有稳定执行 ID，每个已完成步骤保存输出，重启时重新运行控制代码，遇到已有步骤便读取保存值，遇到尚未完成的步骤才执行函数体。程序可以连续表达受理、付款、回填与等待，而执行历史负责跨进程恢复。

DBOS Python 3.3.0 是嵌入应用进程的持久执行库。应用进程执行工作流，系统数据库保存工作流状态、输入、步骤输出及错误。开发可采用把数据库嵌在进程内、以本地文件保存数据的 SQLite，生产部署官方推荐独立服务式的 Postgres。前者启动简单，后者把并发连接、备份和高可用交给数据库部署管理。系统数据库属于新的持久域，它与购买权威可能是不同数据库。基础库使用 MIT 许可，允许按条款使用、修改和分发，并要求在副本或重要部分中保留版权与许可通知。额外的托管或企业控制能力按相应服务确定~\cite{ch21-dbos}。

一个购买编排可以写成三个有序步骤：\code{bind} 读取并核对 $A_J$，返回原 P/C/T、平台范围、键及完整支付输入；\code{pay} 沿这些原值查询或发送；\code{finish} 核验可信结果并提交 $R_J$、O 与 EV。工作流 ID 可以由已保存的完整 S 派生，并与业务绑定建立持久关联。三个步骤的输出分别提供后一个步骤需要的信息，代码中的次序决定重放时如何定位历史。

\CTwentyOneFigure{ch21-13}{已完成步骤从历史恢复，平台效果到步骤检查点之间由原业务身份衔接}{fig:ch21-dbos-replay}

正常路径上，DBOS 先保存工作流输入和状态，再执行 \code{bind}。步骤输出写入系统库以后，后续重放从该输出恢复原绑定。付款步骤向平台发送，得到 X1，再保存步骤输出。最后一步把同一 X1 写回购买权威，并把最终响应作为工作流结果保存。系统库中的步骤索引和序列化输出服务于执行重放，业务权威中的 $A_J/R_J$ 服务于购买仲裁，两份状态各自解释不同对象。

看一个关键窗口。平台余额初始 10000 分，原 C 的 3000 分请求已经提交 X1，应用在保存付款步骤输出之前退出。系统库只有 \code{bind} 的完成输出，付款步骤没有检查点。重启以后，DBOS 跳过已经完成的绑定函数体，再次进入付款步骤。该步骤使用原 \code{scope/pay-C}，平台返回原 X1，余额保持 7000 分，步骤输出随后持久保存并进入 \code{finish}。若步骤每次临时生成付款键，重放会形成新的收款。

步骤完成的判定来自已保存输出。即使函数体在崩溃前已经执行，检查点缺失时它也可能重入。外部服务原键、对象查询和业务权威因此继续负责副作用窗口。DBOS 自身支持的数据库事务集成可以把相应事务与检查点组织在一个支持的边界中，普通 HTTP 支付则位于该边界之外。

工作流控制代码要按相同输入与保存输出产生相同的步骤序列。当前时间、随机数、文件读取、模型回答和外部查询放入步骤，输出能够序列化并保存以后，重放读取原值。若在两个原步骤之间随意插入一个新步骤，旧执行的步骤位置可能与新代码错配。部署为未完成工作流保留兼容代码或应用版本，序列化格式也按历史版本读取。

异常也属于执行结果。配置了步骤重试时，临时失败按有限次数与退避执行；超过次数会向工作流抛出错误，未捕获异常可使工作流进入错误终态。例如平台已提交 X1，步骤却因连接断开报错，此时编排已经结束，原 C/T 继续保存付款未知。恢复负责人沿原键查询，把取得的 X1 回填到权威，再按编排状态决定重启工作流或完成业务收尾。

复用工作流 ID 与比较完整业务输入还有一层区别。调用方以同一个执行 ID 请求原工作流，可以得到已有句柄或结果。假设原输入为 3000 分，后来同 S 携带 3500 分，购买接口先核对 $A_J$ 并返回异参冲突。若只向工作流读取已有结果，可能直接得到原 3000 分结果，缺少对新输入的冲突说明。执行 ID 去重服务于工作流，完整 S 的绑定比较服务于业务入口。

这一结构的成本可直接计量。每个步骤新增输出、索引与状态写入，长流程的历史保留占用数据库；大量同时恢复的工作流争用连接、查询和执行池；大模型输出与工具返回还增加序列化和存储体积。把上千个细碎的纯计算都切成步骤，会放大检查点成本。适合的边界通常围绕非确定性结果、长等待和外部责任安排，使恢复能复用昂贵结果，又保持历史规模可管理。

\subsection{Restate 用独立日志服务承接执行}

Restate 是独立部署的持久执行运行时，把接入、持久日志、状态物化和业务处理端点分开。接入组件 Ingress 接受客户端对服务及其 handler 的调用，handler 就是应用实现的具体处理函数。Ingress 根据调用的工作流身份、对象键或幂等键定位分区，分区日志先保存调用事件。partition processor 是负责该分区的执行者，它读取日志，通过双向流调用应用端点。端点内的 SDK 把持久步骤、状态更新、计时器和子调用转成日志动作，processor 在相应持久记录可读后确认步骤继续。应用因此可以按函数组织业务，执行历史由独立运行时负责保存和恢复~\cite{ch21-restate}。

\CTwentyOneFigure{ch21-22}{数据库式和日志式持久执行分别把历史放在系统库与独立分区运行时，原 C 连接外部支付}{fig:ch21-durable-components}

Restate 的分布式日志层 Bifrost 与执行者分离。每分区的 sequencer 负责给事件分配顺序并复制到日志副本，多数确认定义正常路径提交。processor 把日志应用到嵌入式键值引擎 RocksDB，形成可按键查询的 journal、幂等元数据、对象状态和计时器索引。这里的 journal 是一次调用按步骤保存的执行历史，包含重放所需的动作与结果。跟随者也可持续物化，减少接管后的重建工作。对象存储快照保存某个位置的物化状态，恢复者下载快照后读取其后的日志，降低长历史的回放成本。日志、物化副本和快照各占一份资源，部署可以分别增加日志保存容量与 processor 执行能力。

以 \code{PaymentService/process} 和原 C 为目标，Ingress 先把调用放入对应日志。processor 建立执行并向 handler 发送已有 journal，handler 遇到已完成的 \code{ctx.run} 就读回原输出。新的付款步骤发出原键 HTTP，拿到 X1 后，步骤结果经 processor 写入日志，多数确认并重新读到记录以后，运行时才把该步骤作为已保存事实继续。进程或连接中断时，新尝试取得同一 journal，跳过已有输出。

应用端点按运行时历史重放，随机选择、当前时间和外部 API 放入持久步骤。除了普通服务，Restate 还提供虚拟对象：每个对象键关联持久状态，同一键的独占处理按运行时规则串行推进，只读共享处理使用相应的读取规则。按 C 建对象可以整理一份付款的局部流程，跨账户库存等约束则会越过单键边界。跨不同对象或服务的调用，在来源分区记录后交给内部投递结构，目标分区以其序号仲裁接收。外部银行、仓库或第三方 HTTP 位于这套内部投递域之外，步骤在平台成功后、日志记录前中断时，继续使用原业务身份恢复。

日志切换与执行切换也分开。控制面保存成员、放置、epoch 和日志片段描述。重新配置先封闭旧活动日志片段，确定权威尾部，再发布新的片段与 sequencer。processor 接管取得新 epoch，迟到的旧尝试事件受资格检查。它们保护运行时内部的写入顺序，已离开 handler 的付款请求继续沿原 C 处理。

Restate 服务端 1.8.0 使用商业源码许可 BSL 1.1。它允许读取、修改和再分发源码，生产使用按附加授权判断，并在该版本发布四年后转为 Apache 2.0。该版附加授权允许运行自己组织编写的服务与工作流，限制让第三方直接通过 Restate API 注册并调用其自身服务的公共托管平台。内部购买编排与向外销售这种平台具有不同授权边界。部署除业务端点外，还要为日志节点、processor、元数据多数与快照存储安排故障域和容量。日志副本保存写入，RocksDB 提供低延迟随机访问，快照与日志裁剪要衔接完整恢复点。与 DBOS 的数据库式结构相比，状态和恢复负担移动到独立执行运行时，运维边界也相应增加。

对购买系统，一种接法是让 Restate 保存编排，原分区权威保存 S 与资源仲裁。运行时在 journal 中记录调用和步骤结果，权威在 $A_J/R_J$ 中保存受理与原结果，恢复时由原 C 连接两层历史。另一种接法把键内业务状态放入虚拟对象，相应对象承担业务写入仲裁，跨对象库存通过持久决定连接，退款对象读取原收款并仲裁累计占额。后一种接法改变了事实的保存位置，迁移过程要把旧绑定、资源约束和外部关联一同交接。

DBOS 和 Restate 的例子把付款步骤写进确定的控制代码。模型参与下一步选择以后，执行历史还要固定模型提出了哪项购买、Alice 批准了哪个版本，以及工具重入时应恢复哪个业务身份。

\subsection{Agent 工具调用连接模型、授权与业务责任}

Agent 是由模型提出下一步动作、程序调用工具并把观察交回模型的执行系统。它在模型选择、工具调用、人工确认和长等待之间切换，代码重启以后可能重新走到某次付款工具。LangGraph 用图节点表示计算步骤，边表示继续执行的关系，并通过 checkpointer 保存按运行轮次形成的图状态和调度位置。checkpointer 是可替换的检查点保存接口，选择持久实现可以在进程丢失后恢复，内存保存器只适合进程内运行。thread 是一段逻辑执行历史，其 ID 供恢复时定位原检查点；工具调用 ID 关联模型消息中的某次工具请求，业务 S 则标记购买。图流程可灵活组织分支和人工等待，检查点存储、历史体积与重放兼容也随之进入架构~\cite{ch21-langgraph}。

图内通常包含计划节点、人工批准节点、工具网关和结果节点。计划节点把模型选择规范化成结构化决定，例如购买 \code{paper-001} 1 件、报价 Q7、CNY 3000 分。图检查点保存决定版本 d17、模型输出引用与规范化输入。授权服务让 Alice 批准 d17 中的具体内容，批准记录再与账户、允许动作、金额上限和有效期关联。工具网关读取保存的 S、输入与授权，向购买权威受理。

\CTwentyOneFigure{ch21-17}{模型决定、具体授权和稳定业务身份先保存，工具网关再把调用接到原购买责任}{fig:ch21-agent}

给出完整时序。t0 模型生成 d17，t1 持久保存决定和原 S，t2 Alice 批准 3000 分购买，t3 网关保存授权关联并向购买接口发送，t4 平台完成 X1，t5 工具节点在写回图检查点前崩溃。恢复从已保存的 d17 和 S 继续，调用购买查询得到原 X1 与 O1，再把工具结果写入图。重新运行模型如果产生电子书或 3500 分报价，形成新的决定版本 d18，必须先把原购买结果查明，再按新的业务意图和授权处理。

\CTwentyOneFigure{ch21-28}{人工批准绑定决定版本，业务受理以后以原 S 恢复，模型重规划另建决定与授权分支}{fig:ch21-agent-states}

检查点的保存方式影响可丢失窗口。同步持久模式在推进下一步前等待保存，增加关键路径 I/O；异步模式允许下一步与保存重叠，崩溃时可能重放尚未保存部分；只在退出时保存则形成更大的中间窗口。涉及付款的节点在任何模式下都从持久业务身份进入，检查点丢失可以使工具重入，原购买协议负责汇合外部效果。

人工中断也涉及重执行。LangGraph 的 interrupt 让图保存当前状态并向调用方报告等待输入，外部稍后提供批准或修改，再恢复该 thread。恢复时包含中断的节点可以重新从开头运行，批准之前的副作用因此单独封装并保存身份，或者安排到明确的后续节点。授权服务保存 Alice 对 d17 的批准，网关每次发送前读取该记录，核对账户、动作、金额、有效期与原输入。模型若生成 d18，或者文本中声称用户已经批准，网关根据保存的批准版本比较，版本或内容不符就返回授权冲突。独立批准记录使重新规划、节点重放和用户迟到回复都能定位到具体获批的决定。

当前 LangGraph 的节点超时用于异步节点，达到时限后终止本次等待或触发新尝试。失败尝试的图内缓冲写入可以清除后重试，已经提交到外部平台的付款继续存在。网关对 HTTP 超时保存未知 C/T，恢复按原键查询。若用户决定退款，创建关联 X1 的 RF，成功退款与未知占额一起受 3000 分原收款约束。补偿是新的持久业务动作，其结果也要回到图历史。

长期 Agent 还会积累模型响应、工具输入输出、授权、检查点与等待计时器。模型版本变化、提示变更和代码升级影响继续执行，恢复先读取已经保存的决定，再从相应节点继续。历史清理分别设置图重放、业务去重和审计的保留条件，未知付款关联的原 S、决定与授权保留到相应责任结束；若清理图检查点，则先把恢复所用关联转存到持久业务记录。运维从 thread 找到原 S，从原 S 找到 C、X、EV 和 F，再进入对应责任的查询、回填或补偿接口。

第二个方向沿关键路径寻找等待。保存每个步骤会增加串行持久化，复制命令会经过网络确认，远端读又可能绕到领导者。下面三项研究分别调整这三处路径：libDSE 允许内部状态先推测传播，Jetpack 让适合的命令走并行共识快路径，Bodega 将读取资格放到指定地域。它们各自的确认集合、回滚规则和租约条件在对应协议中成立。

\subsection{libDSE 重叠内部执行与后台持久化}

微服务链每经过一次持久步骤就同步等待保存，短计算、长链路的业务会累积多次持久化往返。libDSE 是把分布式推测执行组织成应用运行时的研究设计：参与者先在可回滚状态中通信，后台持久化继续推进。它采用崩溃后重启的 fail-restart 模型，参与者按协议执行，在有界时间重启，私有持久存储跨崩溃保存，可靠点对点信道在网络隔离消除后继续传输。运行时追踪状态之间的恢复依赖，对外暴露结果以前，等待相关依赖都跨过可恢复边界~\cite{ch21-libdse}。

参与者把应用状态封装为 \code{StateObject}，也就是能按版本保存和恢复的状态对象。应用实现四个存储接口：\code{Persist} 保存指定版本，\code{Restore} 恢复已有版本，\code{Prune} 删除已经可回收的旧版本，\code{ListVersions} 枚举实际可恢复的版本。版本不仅包含业务状态，还带运行时的依赖元数据。action 是应用的一段原子状态访问，它与保存、恢复互斥，使这段访问的效果整体进入某版本或在回滚中整体丢失。不同 action 之间的应用并发保护由应用提供，余额等共享资源按自己的锁或事务规则仲裁。封装版本把恢复交给运行时，同时要求应用提供正确的快照、恢复和并发边界。

长 RPC 等待放进 \code{sthread}。它是一段脱离当前 action、携带推测依赖集合的执行，可以等待、重试或发消息。返回结果重新合并到父 StateObject 时，运行时检查产生该执行的状态是否已经回滚。若已回滚，合并失败，该分支停止，父对象按恢复后的状态重新推进。长等待因此不必把 StateObject 的持久化全部阻塞。对外输出前调用 \code{Barrier}，这是等待相关状态及其依赖全部可恢复的屏障。内部调用可以先利用推测结果，屏障把等待集中到真正跨出内部恢复域的位置。

\CTwentyOneFigure{ch21-14}{内部推测状态可以先传播，输出跨过 Barrier 以后才暴露给外部实体}{fig:ch21-dse-boundary}

依赖图的顶点由对象 ID、全局故障序号和本地持久化计数标识。顶点 $O^7_4$ 表示订单对象在故障序号 7 下、本地持久化计数为 4 的状态点。这个顶点可能还在后台保存，是否可恢复另由存储完成情况判断。边 $u\rightarrow v$ 表示恢复 u 却没有 v 会产生不一致：同一对象的后续版本依赖前缀，消费消息的对象依赖消息来源版本。控制流程也在图中，订单因为库存确认而发起付款，即使没有直接读库存内存，也依赖那份确认。

假设库存对象 I 在版本 $I^7_4$ 暂时预留 SKU-A，向订单对象 O 返回推测确认。O 在 $O^7_4$ 保存“可以付款”，其依赖集合包含 $I^7_4$。O 随后把结果交给执行分支 W，W 在支付前调用 \code{Barrier}。此时 O 的文件先保存完成、I 尚未完成，W 继续等待。若 I 随后完成持久化，协调器可找到包含两者的闭包，W 得到继续资格。若 I 先崩溃并只能恢复到版本 3，O 中依赖版本 4 的状态随协议回滚，W 被放弃。

\CTwentyOneFigure{ch21-23}{依赖边连接订单控制与库存确认，持久闭包允许输出，丢失顶点引发相关回滚}{fig:ch21-dse-graph}

可恢复边界是一个闭包：其中顶点全部可恢复，所有依赖边也落在闭包内。前例中，即使 $O^7_4$ 的文件已经保存，它所依赖的 $I^7_4$ 也要保存完成，二者才进入同一闭包。协调器接收参与者持久保存的图片段，在内存图上搜索闭包并发布边界。协调器重启后从这些持久片段重新计算正常路径的边界，因此每次边界推进可以省去单独同步保存全局决定的写入。

持续通信可能让依赖追向未来，出现始终找不到有限恢复范围的多米诺效应。libDSE 限制本地计数为 y 的状态接收来自计数 n 的消息时满足 $y\ge n$。若本地计数落后，运行时发起一次持久化以推进本地计数，然后可以开始 action，无须等实际持久写完成。参与者频繁通信时会被动地接近相似的持久化进度，有限计数前缀因而能形成闭包。

故障恢复要对回滚范围形成一致决定。协调器根据故障对象报告的持久版本删除丢失顶点，再反复删除指向丢失部分的依赖顶点，直到剩余图闭合。它把新的故障序号和回滚动作同步保存到持久日志以后，才发布给参与者。参与者按自己的速度恢复，消息只在相同故障序号间接收。来自旧序号的消息被丢弃，来自更新序号的消息等待本对象跟上。

协调器重启先回放成员与回滚日志，再收集各参与者的持久图片段，完整重建视图后才回答当前边界查询。运行时的 Connect 接口把合法实例注册到当前成员与故障序号中，后台 Refresh 周期交换持久进度和依赖片段，使协调器更新可恢复边界。应用存储层把写入资格授予当前实例，例如用独占卷限制挂载者，或由资源端核验代次并拒绝过时写入。旧实例即使迟到恢复，其消息受故障序号检查，其持久写入受资源端资格检查。

记录粒度取状态点而非单条消息，可以降低正常路径开销，也可能把同一点内互不相关的修改一起回滚。性能代价包括消息元数据、版本存储、图扫描、后台协调及故障时重做工作。收益主要来自原先由串行持久化主导的链路，长模型推理占据大部分时间时，可重叠的持久等待只占较小比例。进展要求本地保存和恢复最终完成，后台进度刷新持续运行，并有足够长的稳定窗口让相应闭包形成。

W 跨过 Barrier 以后，订单和库存的相关依据已经可恢复，随后发出原 C 的 HTTP。平台提交 X1 到本地记录结果之间若发生故障，恢复者通过原键和查询取得 X1，接续付款记录。内部 StateObject 的回滚由依赖图处理，平台收款由外部协议连接。这项设计适合能够封装状态、保存版本并重做内部计算的服务链，输出等待集中在 Barrier，昂贵的内部计算与后台保存则可以重叠。

\subsection{Jetpack 给原共识协议增加并行快路径}

跨地域复制的另一项成本是串行通信轮次。Jetpack 为独立提交的命令增加快路径，客户端把命令同时送到快路径与原协议，原路径负责排序和提交，快路径在没有冲突并取得足够确认时先完成。协议采用 $2f+1$ 个按协议执行的副本，容忍 f 个崩溃故障。副本在两条路径之间协调顺序承诺，使适合的请求减少一次串行往返~\cite{ch21-jetpack}。

每个副本都维护原日志和快路径记录。快路径的视图标记一轮可以共同作出顺序承诺的成员与提议状态，换主时通过恢复切换，避免混用旧新承诺。收到快命令 A 后，副本检查它与两条路径的未提交命令是否冲突，视图一致且无冲突时记录 A 并确认。原路径提议者是能在原协议中安排命令顺序的角色，它还承诺：在最终提交 A 之前，不会把并发冲突命令排到 A 前面。客户端的快确认集合至少包含
\[
q_{\mathrm{fast}}=f+\lceil f/2\rceil+1
\]
个副本，并覆盖原路径全部提议者。三副本时 $f=1$，普通多数为 2，快路径需要 3；五副本时 $f=2$，快路径需要 4。

\CTwentyOneFigure{ch21-15}{两条路径可正常完成命令，视图变化才进入单独的恢复流程}{fig:ch21-jetpack}

以两条已经固定各自 P/C/T、账户和库存均独立的命令 A、B 为例，应用比较交换次序以后的完整状态与各自响应，证明可交换后，快路径可以分别取得确认。首次受理若共享 b1 的编号分配器，交换次序会改变哪个 S 得到 \code{C:b1:1}，账户不同也要判为冲突。冲突集合还包含共享会话额度、Outbox 序号和完整结果字段。若 A 与 B 都从 Alice 的余额扣款，或者都竞争最后一件 SKU-A，原路径同样决定顺序，后执行的命令读取前一命令已经改变的资源状态。

快确认的收益依赖地理位置。三个副本中，客户端要等最远的 C，A--C 的长链路或 C 的慢磁盘可能让原路径更早完成。客户端取先完成的有效路径，饱和时可减少快路径使用比例，以降低同时向多副本发送的扇出与冲突检查开销。满足快提交条件的单条命令经过一次网络往返，即 1 RTT；购买受理命令返回以后，协调者接着完成库存决定，付款工作者调用平台，再由权威提交结果。低冲突、跨域往返占主要成本且快确认集合延迟可接受的命令，更容易从这条附加路径获益。

最困难的部分在换主。A 已经快提交给客户端，原路径还没提交 A，原提议者随后退出。新提议者如果只保留自己已有的原日志，可能把旧未提交冲突命令 B 放到 A 前面。快路径对外已经给出 A 的顺序承诺，新任期必须找回 A，并在冲突命令前落实该承诺。

Jetpack 为快路径单独维护视图，只有两条路径位于同一视图的确认才能组合成快提交。切换时，恢复协调者先确定最后一个正常视图，用 BEGINRECOVERY 让多数暂停该视图的快请求。随后用标准 Paxos 实例仲裁恢复集合：PREPARE 收集旧快日志及已经接受的集合，ACCEPT 让多数保存同一选择。把 Paxos 的轮次规则展开，若回复带有已经接受的集合，协调者选择其中接受轮次最高的值；全部回复均无已接受值时，才按出现次数从旧快日志计算候选集合。新的恢复尝试接续先前可能已经形成的选择，再把确定的集合整体交给原协议提交。

当收集 $f+1$ 个回复时，一个曾得到快多数的命令至少应在其中留下 $\lceil f/2\rceil+1$ 份记录，这给恢复筛选提供依据。若两个冲突命令都达到候选阈值，协议可推出这两个命令均未快提交，选择其中一个加入恢复集合是安全的。恢复集合或空集合对应的 no-op 在原路径提交后，形成新视图的稳定标记。no-op 是不修改业务状态的空命令，在这里承担记录恢复边界的作用。后继换主从标记识别已完成恢复的边界，避免把未完成恢复的中间视图误当作正常视图。

\CTwentyOneFigure{ch21-24}{新视图先确定同一恢复集合，再由原协议提交恢复集合和标记，随后恢复普通提议}{fig:ch21-jetpack-recovery}

原协议接纳快路径时，提议者按快确认承诺安排冲突命令，执行者按最终确定的顺序应用。Jetpack 的公开研究实现把快路径接入 Raft 等不同原协议。接入购买系统以后，复制层增加快日志、客户端扇出、冲突索引与恢复流程；业务层根据完整 S 绑定原输入，用工作代次核验回填，再由跨组 TX 连接多个库存组。复制层确定受理命令的位置，业务层把该命令产生的责任接到后续付款链。

\subsection{Bodega 将线性化读取移到选定地域}

地域入口既有写入，也有大量原结果和状态查询。经典领导者租约让其他副本在约定期限内承诺不支持冲突领导者，使持有有效资格的主副本可以本地回答线性化读，远端入口往往还要访问主副本。Bodega 是把这种读取资格分配到指定地域的研究系统。它用 roster 保存当前领导者及各键范围的 responder 集合，responder 就是获准回答该范围本地读的副本。集群通过租约保持名单一致，写入提交的确认集合则覆盖该键全部 responder，使它们在任何相关写入完成前获得对应信息。降低远端读取延迟的代价，由名单中的写确认与资格维护承担~\cite{ch21-bodega}。

设五个副本 S0--S4，S0 为领导者，某键的 responder 为 S0、S3、S4。S4 持有稳定 roster 并追上安全位置后，就近客户端可以从 S4 读取。写入先发到多数且覆盖 S3、S4。S4 如果收到该键最新槽位且已知提交，立即回答；若该槽位尚未确定，读挂在该槽位的等待集合中，提交通知到达后释放。连续写入不会让这次读无限追逐未来，它等待自己关联的那一个槽位。

\CTwentyOneFigure{ch21-25}{名单指定地域读响应者，写确认覆盖响应者，未确定槽位上的读取等待相应提交}{fig:ch21-bodega}

可选的提前接受通知让副本把写接受信息发给 responder。responder 收集到足够证据，便能更早确认该槽位必然提交，减少读等待。它增加写侧通知流量，写入繁忙时可以关闭。客户端还可以设置 unhold 期限，即允许本次读在该 responder 上等待多久。超时后用同一读请求 ID 向另一 responder 或领导者发送，使用最先到达的有效结果。较短期限更快绕过慢副本，也会增加重复读流量。

本地读资格来自节点两两交互的 roster leases。每份授予绑定同一 ballot，ballot 标识当前名单与领导选择所属的一轮，接收者据此把同一轮的同意组成多数。responder 持有多数有效授予后，继续把日志提交到协议计算的安全阈值，覆盖授予者报告的相关已接受前缀，也就是从日志开始到相应位置的连续历史。持有租约的旧副本在追赶阶段等待，达到该阈值以后才开放本地权威读。

租约依赖有界时钟速率漂移，不要求各节点墙钟时间戳同步。授予者按较保守的期限等待，持有者提前认为租约到期，保证旧名单的读取资格先结束，新名单才生效。切换名单先撤销旧租约，无法联系旧节点时等安全过期，再建立新租约与前缀。节点进程暂停、时钟漂移和虚拟化环境的计时行为因此进入部署前提。

Bodega 把读延迟成本移向 responder 覆盖、心跳租约、每键范围名单和读等待结构。名单包含远端慢节点时，写提交就受其影响，故障后要调整名单并完成租约交接。对不可覆盖的已完成 $R(S)$，业务可以按前述规则复用结果缓存；对尚在改变的责任、会话关闭或准入状态，部署可在读集中的地域配置 responder，并把写入必须覆盖这些地域的延迟纳入预算。论文在 Summerset 中提供研究实现，名单大小与地域位置是组合读写成本的直接参数。

第三个方向把持久状态放到独立对象存储，计算进程可以替换，恢复者按日志和元数据找回有效状态。这里又有两种粒度：SlateDB 组织逐键修改、WAL 和有序文件，Iceberg 与 Delta Lake 组织表级文件集合。先读哪份元数据、何时确认写入持久、哪次提交使候选文件可见，决定了每种结构的恢复步骤。

\subsection{SlateDB 将 LSM 的持久层移到对象存储}

对象存储提供独立于计算进程的持久容量，API 延迟和请求费用通常高于本地磁盘。SlateDB 0.17.0 是采用 LSM 组织数据的嵌入式键值数据库，把日志及有序表文件放入对象存储。LSM 先把修改写入内存有序表 MemTable，随后刷成不可变有序文件 SST，后台整理合并不同层的文件与重复键版本。读取结合内存和文件索引找到最新可见值。计算进程还可缓存读过的数据块及完整 SST，用本地空间换取较少的远端读取。顺序批量写有利于吞吐，合并会再次读写数据，查询也可能检查多个层次，因而分别产生写放大和读放大~\cite{ch21-slatedb}。

在 SST 形成以前，WAL 保存足以重建修改的顺序日志，计算进程丢失后从它重放。SlateDB 把这份日志也写入对象存储，并用 manifest 元数据描述有效日志、文件集合、序号和写者代次。数据文件先写出，元数据再发布有效集合，读者和恢复者据此判断哪些对象属于当前数据库。对象本身的持久保存、有效集合发布和应用得到持久确认分别形成可检查的边界。

\CTwentyOneFigure{ch21-26}{写入句柄、对象日志、SST 与 manifest 分别位于内存受理、持久化和有效文件集合的边界}{fig:ch21-slatedb}

一份写入首先更新内存 WAL 和 MemTable，API 返回 \code{WriteHandle}。这是关联本次写入完成状态的句柄，调用者通过 \mbox{\code{await\_durable}} 等待它达到持久完成，或者调用 flush 促使并等待相关批次刷出。把写入与等待分开，调用者可以先积累多份写入，再共同等待持久边界。内存中已经可见和跨计算进程故障可恢复是两个时刻。若把候选付款结果缓存写入 SlateDB，业务在承诺保存之前等待所需持久边界。

取启用对象 WAL 并等待 \mbox{\code{await\_durable}} 的路径：记录 \code{C:b1:1 -> X1} 在时间 t0 写入 MemTable，t1 返回句柄，t2 对象 WAL 保存完成并使等待者获知持久完成，t3 后台把 MemTable 刷成 SST。若进程在 t1 后、t2 前退出，该写可能没有成为可恢复事实。t2 以后恢复者从持久日志重建这份修改，t3 形成的 SST 则缩短后续重放并改变读取成本。

批量把大量小修改摊到少量 PUT，减少调用费用，却增加等待和内存规模。设每秒均匀到达 1000 个 1 KiB 写入，每份单独 PUT 对应每秒约 1000 次调用。若按固定 100 毫秒形成一批，每秒约 10 批、每批约 100 KiB，再加日志与文件元数据。写入在这 100 毫秒内均匀到达时，纯等批造成的平均附加等待约 50 毫秒，持久确认还要加上对象 API 时间。缩短批次周期会降低等待、增加请求次数，延长周期则占用更多缓冲内存。

manifest 中的 \code{writer\_epoch} 标记当前写者代次。新写者提高代次，旧写者观察到更高代次后返回 \code{Fenced} 错误，终止过时写入。计算位置可以改变，有效文件集合与写入资格通过对象上的元数据协议衔接。后台 compaction 把旧有序文件整理为候选新文件，发布新集合以后，旧文件按检查点、读者与保留规则回收。大量小文件会增加读取与元数据开销，积极整理则消耗对象请求、网络和计算，整理速率随写入量与恢复预算选择。

S3 支持 \code{If-None-Match} 条件创建和基于 ETag 的 \code{If-Match} 条件更新。ETag 是服务返回的对象实体标识，反映对象内容变化，条件请求用它比较本次读取与当前对象是否匹配，独立完整性检验则使用应用选择的内容摘要。设上层发布协议在 manifest 内容中保存发布代次 generation，由发布者每次提交时增加一代。两个写者都读到 generation 7 与 ETag \code{e\_old}，第一个用带 \code{If-Match: e\_old} 的 PutObject 发布 generation 8 的新内容，成功后服务返回 \code{e\_new}。第二个再按 \code{e\_old} 提交时得到 412 响应，读取当前内容与 ETag 后重新构造提交。并发删除还可能产生冲突 409 或未发现 404，恢复者按具体操作重新判断对象和候选上传，分段上传的某些失败要求重新启动上传。版本化桶的当前删除标记也会改变条件创建的判断。上层协议由此先保存数据对象，再以条件写仲裁新集合，读者固定发布结果后读取相应数据~\cite{ch21-s3-conditions}。

SlateDB 的部署把对象存储凭据、条件请求权限、WAL 与文件保留、缓存冷启动、整理带宽和读放大纳入日常运维。升级先核对格式与 API 兼容规则，让新代码读取已有 manifest、WAL 与 SST，再切换写者。对购买系统，大规模执行状态、冷索引和派生视图可以利用这种计算与保存分离的结构；若业务权威采用它作为底层引擎，状态机在其上保存完整 S 与结果，并把同一命令涉及的资源写入、外部责任和跨组决定组织成可恢复协议。

\subsection{Iceberg 和 Delta 发布一组文件的可见状态}

湖仓表把对象存储中的数据文件组织成可查询、可事务更新的表，多种计算引擎通过共同表格式读写。Apache Iceberg 与 Delta Lake 都属于这种表格式及相应实现，它们管理的是一个一致的文件集合。工作者先写候选数据文件，再用表级提交发布其中一组。读者从表快照或事务日志取得有效文件集合，按引用读取数据，尚未发布的候选文件等待提交，失败尝试留下的孤立文件按保留规则回收。计算与保存能够分别扩展，表元数据、提交服务与引擎兼容则成为新的管理对象。

Iceberg 1.12.0 把表状态组织成元数据树。表元数据指向快照，快照代表一次完整可见的文件集合；manifest 列出数据及删除文件和相应状态，manifest list 再列出这次快照引用的各份 manifest。分层索引使更新可以复用未改变的子树，避免每次复制全部文件清单。目录服务 Catalog 保存当前表元数据的位置，并通过它所支持的原子提交把指针由旧文件换到新文件。读者先固定快照，再按其引用读取，长期读取不会因为后台出现新候选文件而改变本次集合。目录的条件更新与可用性因而成为表提交的关键依赖~\cite{ch21-iceberg}。

\CTwentyOneFigure{ch21-16}{候选数据与表级提交分开，重试验证当前假设以后发布新的完整集合}{fig:ch21-lake}

设当前元数据 M7 指向快照 S7，包含文件 A、B。导出作业 J1 为来源截止 L9000 写出 C，计划追加 C。整理作业 J2 同时把 A、B 合成 D。J2 先提交 M8，文件集合变为 D。J1 的 M7 条件提交失败后读取 M8，验证追加假设，复用已经生成的 C 和对应 manifest，再提交 M9，读者看到 D、C。候选文件写入没有全部重做，提交顺序却保持为 M7、M8、M9。

若 J1 的操作也是重写 A、B，J2 已经删除其有效引用，J1 就必须重新验证或重新计算。乐观提交只是在最后竞争一个发布位置，操作能否复用取决于先前读取的假设是否继续成立。整理需要原文件依然有效，追加往往更容易复用元数据。目录的原子比较交换、事务目录或支持的提交实现，构成这段协议真正的仲裁点。

Delta Lake 4.4.1 用递增版本的事务日志组织表状态，每次提交记录一组动作，例如添加文件、删除文件、改变元数据或协议状态。读者从日志及其检查点恢复某个版本的有效文件集合，形成版本化快照。写者读取当前版本，写候选文件，验证从读取版本以来的并发改变，再提交新版本。日志把一批文件变化和相关事务标记放在一次发布中，长历史则要依靠检查点和日志保留控制重建成本。实际日志发布依赖所选存储或目录实现提供正确的互斥与提交条件~\cite{ch21-delta}。

\CTwentyOneFigure{ch21-27}{Iceberg 用元数据树和目录指针发布，Delta 用版本日志描述文件动作，来源批次身份贯穿两条路径}{fig:ch21-lake-structures}

例如两个 MERGE 作业分别处理 10 月 17 日和 18 日数据。MERGE 根据匹配条件更新已有行或插入新行，因此先读取目标表，再发布相应文件动作。若谓词只关联用户与来源日期，而目标读取范围没有明确限定日期，查询规划可能扫描两天文件，其中一个作业新增或删除文件便会与另一个的读取冲突。把目标日期和其他分区条件明确写入操作，可以使各作业读取相应分区的文件。遇到表级元数据或协议版本改变时，写者另外验证自己使用的模式和提交规则，再按新版本重新规划。

表事务与来源去重分别建立。若账本快照 L9000 导出重试两次，每次都成功追加相同业务行，表级串行提交可以得到两个合法版本，却含重复来源。导出作业保存批次 ID、来源分区和截止、数据摘要、目标表及提交版本，重启先判断原批次是否已经发布。

Delta 的 \code{txn} 动作能把应用 ID 与应用版本同文件动作一同提交。设导出流 \code{ledger-b1} 的第 42 批读取固定来源，写出文件 C，同时提交 \code{appId=ledger-b1, version=42}。应用重启先读取快照中该 appId 的版本，若按自身单调批次规则确认 42 已完成，就跳过追加。版本 41 则重新尝试第 42 批，文件动作与版本动作一起提交。Delta 保存每个 appId 的最新版本，单调递增和批次输入不变由应用定义。保留与并发运行也遵循所选引擎规则，避免两个独立作业错误共用身份。Iceberg 的相应导出方案由作业和目录保存批次提交关联，读回原提交后再决定继续。

选择表格式时，先核对使用的查询和写入引擎、Catalog 或日志提交实现，以及重写、追加和模式变更的并发规则。大量文件的追加与长期快照读取会放大元数据和保留成本，频繁小批更新则产生较多候选文件、整理和冲突重试。把批次做大可摊薄提交开销，也会增加数据到达查询者的等待。导出延迟、历史审计和对象请求费用共同决定批量与保留策略。

对账表回滚到旧快照时，分析读者重新看到旧文件集合，平台保存的 X1 与购买权威的 $R(S)$ 继续描述原收款。批处理发现坏数据时，先纠正导出与表状态，确有业务退款需求时再创建 RF，由退款协议处理资金返还。快照过期与文件清理根据已提交引用和长读者回收对象，并为重跑与审计保留相应时间窗口。
